% SPDX-License-Identifier: CC-BY-SA-4.0
% Copyright (c) 2026 Rafael Beraldo
% Parte 5 (continuação): o rádio portátil, antenas, programação e dicas
% Incluído por presentation.tex; compile a partir dele.

% --- O rádio portátil (HT): anotações em figuras/baofeng.tex, uma página por slide
\begin{frame}{Conhecendo os equipamentos básicos}
  \begin{columns}[c,onlytextwidth]
    \column{0.42\textwidth}
      \centering
      \includegraphics[height=6cm]{baofeng-vfo}
    \column{0.55\textwidth}
      \textbf{O rádio portátil (HT)}, aqui um Baofeng UV-5R
      \begin{itemize}
      \item \alert{Dual-band}: 2\,m (VHF) e 70\,cm (UHF)
      \item Peso: cerca de 210\,g, com bateria e antena
      \item Preço: cerca de R\$ 100
      \end{itemize}
  \end{columns}
\end{frame}

\begin{frame}{Anatomia do rádio}
  \begin{columns}[c,onlytextwidth]
    \column{0.52\textwidth}
      \figura[page=1]{baofeng}
    \column{0.46\textwidth}
      \small
      \begin{enumerate}
      \item \textbf{Visor}: mostra a frequência ou o canal
      \item \textbf{Knob}: liga, desliga e ajusta o volume
      \end{enumerate}
  \end{columns}
\end{frame}

\begin{frame}{Duas linhas do visor: A e B}
  \begin{columns}[c,onlytextwidth]
    \column{0.52\textwidth}
      \figura[page=2]{baofeng}
    \column{0.46\textwidth}
      \small
      \begin{itemize}
      \item O rádio está \alert{sempre ouvindo} as duas linhas
      \item O botão \textbf{A/B} escolhe por qual linha você transmite
      \item A seta no visor mostra a linha ativa
      \end{itemize}
  \end{columns}
\end{frame}

\begin{frame}{PTT: aperte para falar}
  \begin{columns}[c,onlytextwidth]
    \column{0.52\textwidth}
    \figura[page=3]{baofeng}
    \column{0.46\textwidth}
    \small
    \begin{itemize}
    \item Aperte, espere um segundo, fale a um palmo do rádio e solte
    \item Enquanto você fala, você \alert{não ouve} ninguém
    \item O rádio fala com \alert{todos ao seu alcance}: não há privacidade\asterisco % a criptografia é crime
    \end{itemize}

    \bigskip\raggedleft
    \comentario{O subtom controla apenas o squelch e não torna uma transmissão privada; além disso, a criptografia é crime.}
  \end{columns}
\end{frame}

\begin{frame}{Memória: VFO e MR}
  \begin{columns}[c,onlytextwidth]
    \column{0.52\textwidth}
      \figura[page=4]{baofeng}
    \column{0.46\textwidth}
      \small
      \begin{itemize}
      \item \textbf{VFO}: você digita a frequência
      \item \textbf{MR}: canais salvos na memória, com nome
      \item O botão \textbf{VFO/MR} alterna entre os dois
      \end{itemize}
  \end{columns}
\end{frame}

\begin{frame}{Antenas e programação}
  \begin{columns}[T,onlytextwidth]
    \column{0.48\textwidth}
      \textbf{Antenas}
      \begin{itemize}
      \item Primeiro upgrade; a “rubber ducky” original é um elo fraco
      \item J-Pole, dobrável, direcional (yagi ou log-periódica)
      \item Transmissão na vertical e o mais alto possível
      \item Nunca transmita sem antenas
      \end{itemize}
    \column{0.48\textwidth}
      \textbf{Programação com o CHIRP}
      \begin{itemize}
      \item Gratuito e aberto: \linkexterno{\linkCHIRP}{chirpmyradio.com}
      \item Programação de canais e de opções do rádio
      \item Sincronia entre rádios do mesmo grupo
      \item Cabo USB com conector K1
      \end{itemize}
  \end{columns}
\end{frame}

\begin{frame}{Algumas dicas antes de sair de casa}
  \small
  \begin{dica}
    Transporte o rádio \alert{sem a bateria encaixada}: o knob gira na mochila, o rádio liga e
    chega descarregado. Já aconteceu.
  \end{dica}

  \pause
  \begin{dica}
    Fique atento ao \alert{squelch} (SQL): silencia o chiado quando ninguém fala; alto demais, corta sinais fracos. O subtom (CTCSS) \alert{não dá privacidade}: só faz o squelch abrir quando o sinal traz aquele tom.
  \end{dica}

  \pause
  \begin{regulamento}
    Rádios como o UV-5R transmitem nas faixas de radioamador: usá-los exige \alert{COER e licença} com indicativo. Eles não são rádios de uso geral (RUG).
  \end{regulamento}
\end{frame}

%%% Local Variables:
%%% TeX-master: "../presentation"
%%% fill-column: 100
%%% End:
